\section{Utilidad del modelado del subsuelo, métodos disponibles y brecha de acceso}
\label{sec:problema-modelado}

\subsection{El problema no es abstracto: se toman decisiones sobre un medio que no se ve}

Toda obra se apoya en un material cuya estructura no puede observarse directamente en
su totalidad. El suelo no es una pieza fabricada y uniforme: su composición, rigidez,
compactación, saturación y espesor cambian con la profundidad y también entre puntos
cercanos. Por eso, antes de construir, rehabilitar o estudiar una obra aparecen
preguntas concretas:

\begin{itemize}
    \item ¿A qué profundidad cambia significativamente la rigidez del terreno?
    \item ¿Existe una capa blanda, un relleno, una cavidad o un contacto con roca?
    \item ¿El perfil observado en un sondeo representa realmente todo el sitio?
    \item ¿Cómo puede amplificar el terreno una vibración o un movimiento sísmico?
    \item ¿Dónde conviene realizar ensayos adicionales antes de tomar una decisión?
\end{itemize}

Modelar el subsuelo significa transformar observaciones parciales en una
representación de capas o zonas con propiedades mecánicas estimadas. El modelo no es
una fotografía exacta del terreno; es una hipótesis cuantitativa que debe ser compatible
con las mediciones, la geología conocida y la resolución del método empleado.

En este trabajo interesa especialmente la velocidad de onda de corte, $V_S$, porque se
relaciona con el módulo de corte a pequeñas deformaciones mediante

\begin{equation}
    G_{\max}=\rho V_S^2,
\end{equation}

donde $\rho$ es la densidad del material. Un perfil $V_S(z)$ permite reconocer cambios
de rigidez con la profundidad y aporta una entrada fundamental para la clasificación
sísmica y el análisis de respuesta de sitio~\cite{Kramer1996,Foti2014}. Cuando se
repiten perfiles a lo largo de una línea o sobre una grilla, la información deja de ser
puramente puntual y puede convertirse en una sección 2D o en un modelo 3D.

La utilidad no está, entonces, en ``medir una señal sísmica''. La señal es el dato
crudo; el producto de ingeniería es un modelo del subsuelo con alcance, supuestos e
incertidumbre declarados.

\subsection{Casos reales: de la medición a una decisión}

Los siguientes casos muestran que el modelado sísmico del subsuelo ya se utiliza para
problemas de construcción, rehabilitación, clasificación de sitio e investigación. No
son demostraciones de laboratorio: en todos ellos la información obtenida responde a
una necesidad concreta del sitio.

\subsubsection{CentrePort, Wellington: comprender daños y orientar la mitigación}

El puerto de Wellington, Nueva Zelanda, sufrió licuefacción, desplazamiento lateral y
daños estructurales durante el terremoto de Kaik\={o}ura de magnitud $M_w=7{,}8$. Para
comprender la respuesta del terreno y proponer medidas de mitigación se combinaron
\textbf{114 mediciones HVSR}, ensayos MASW activos y arreglos pasivos MAM. En seis
ubicaciones se desarrollaron perfiles $V_S$ de más de \SI{200}{\meter} de profundidad,
con los cuales se estimaron la rigidez, la profundidad a roca y el período fundamental
del sitio~\cite{Vantassel2018}.

El resultado no fue un número aislado. Se obtuvo un mapa que reveló cambios abruptos de
la estructura del subsuelo en distancias cortas. Además, los períodos estimados con
ruido ambiental fueron coherentes con los períodos de máxima amplificación registrados
durante terremotos reales. Este caso muestra dos ideas importantes para la tesis: la
información espacial puede cambiar una interpretación de sitio y la combinación de
métodos activos, pasivos y datos de sondeo es más fuerte que cualquiera de ellos por
separado.

\subsubsection{Mina abandonada en Colorado: localizar zonas peligrosas antes de construir}

En una antigua mina de oro de Colorado se necesitaba evaluar el terreno antes de
construir un nuevo estanque industrial. El problema eran los pozos colapsados, rellenos
y labores mineras someras que habían quedado pobremente documentadas. Se adquirieron
\textbf{40 perfiles MASW} a lo largo de varios trazados para investigar aproximadamente
los primeros \SIrange{9}{15}{\meter} del subsuelo. A partir de las anomalías detectadas
se elaboró un mapa de posibles zonas débiles o vacíos para orientar los trabajos de
movimiento de suelo y construcción~\cite{Crook2017}.

El valor del método fue ampliar la cobertura entre observaciones puntuales. El estudio
no ``demostró'' por sí solo la geometría exacta de cada labor minera, pero permitió
concentrar la investigación y la gestión del riesgo donde la respuesta del terreno era
anómala.

\subsubsection{Delhi Technological University: clasificar varios sectores con un equipo comercial}

En el campus de Delhi Technological University se realizaron 30 ensayos MASW en seis
sectores. La adquisición utilizó un sismógrafo PASI GEA24 de 24 canales, geófonos
verticales de \SI{4.5}{\hertz}, espaciamiento de \SI{2}{\meter} y una maza como fuente.
El trabajo reportó una velocidad de corte media de \SI{285}{\meter\per\second} y
clasificó el sitio dentro de la clase D utilizada por NEHRP~\cite{Gautam2020}.

Este caso es especialmente útil para esta presentación porque conecta una aplicación
real con uno de los equipos comerciales analizados más adelante. Muestra el flujo
industrial completo: arreglo multicanal, fuente controlada, adquisición, curva de
dispersión, inversión y clasificación del terreno.

\subsubsection{Metro Manila: MASW y SPT como fuentes complementarias}

En tres sitios de Metro Manila se comparó la estimación de $V_{S,30}$ obtenida mediante
investigación geotécnica ---valores SPT y calidad de roca--- con la obtenida mediante
MASW. Aunque los valores numéricos no fueron idénticos, ambos caminos produjeron una
clasificación sísmica NEHRP consistente en los tres sitios~\cite{Monjardin2023}.

La lección no es que MASW reemplace al SPT. El sondeo aporta muestras y una descripción
local detallada; el método geofísico aporta cobertura y una estimación dinámica. Su
comparación permite detectar contradicciones, reducir supuestos y construir una
interpretación más defendible.

\begin{table}[ht]
\centering
\small
\caption{Qué demuestra cada caso de estudio.}
\begin{tabularx}{\textwidth}{@{}>{\raggedright\arraybackslash}p{0.19\textwidth} Y Y@{}}
\toprule
\textbf{Caso} & \textbf{Pregunta práctica} & \textbf{Resultado utilizable} \\
\midrule
CentrePort & ¿Por qué varió la amplificación y qué profundidad debe incluir el modelo? & Mapa de período de sitio y perfiles profundos de rigidez y roca para análisis de respuesta. \\
Mina de Colorado & ¿Dónde pueden existir labores colapsadas antes de construir? & Mapa de zonas débiles o vacíos para orientar investigación y movimiento de suelo. \\
Campus de Delhi & ¿Cómo varía la rigidez entre sectores del campus? & Perfiles $V_S$ y clasificación sísmica a partir de 30 ensayos. \\
Metro Manila & ¿Son compatibles la investigación geotécnica y la geofísica? & Clasificación NEHRP consistente mediante SPT/RQD y MASW. \\
\bottomrule
\end{tabularx}
\end{table}

\subsection{Métodos disponibles: cada uno observa una parte distinta del problema}

No existe un método universalmente superior. La elección depende de la propiedad que
se desea conocer, la profundidad, la resolución, la cobertura, el acceso al sitio y el
nivel de incertidumbre aceptable. Una primera división separa los ensayos que requieren
penetrar o perforar el terreno de aquellos que infieren propiedades desde la superficie.

\begin{table}[ht]
\centering
\footnotesize
\caption{Comparación funcional de métodos de caracterización.}
\begin{tabularx}{\textwidth}{@{}>{\raggedright\arraybackslash}p{0.13\textwidth} >{\raggedright\arraybackslash}p{0.18\textwidth} Y Y@{}}
\toprule
\textbf{Método} & \textbf{Información principal} & \textbf{Fortaleza} & \textbf{Limitación relevante} \\
\midrule
SPT & Resistencia a penetración y muestra alterada & Práctica extendida, contacto directo y correlaciones geotécnicas conocidas & Información puntual; requiere sondeo; $V_S$ se obtiene sólo mediante correlaciones empíricas \\
CPT/CPTu & Resistencia de punta, fricción y, con CPTu, presión de poros & Perfil casi continuo y buena resolución estratigráfica & Intrusivo; no recupera muestra; dificultad en gravas o materiales muy duros \\
Downhole / crosshole & Tiempos de viaje y $V_P$/$V_S$ en profundidad & Alta resolución local y referencia valiosa para validar otros métodos & Requiere uno o varios pozos y representa un volumen limitado \\
Refracción / tomografía & Velocidad de primeras llegadas, usualmente $V_P$ & Buena lectura lateral de contactos y profundidad a roca & Las capas de baja velocidad bajo capas rígidas pueden quedar ocultas; no entrega directamente $V_S$ salvo configuración específica \\
MASW activo & Curva de dispersión y perfil $V_S(z)$ & No invasivo, multicanal y repetible; permite perfiles 1D y secciones pseudo-2D & Requiere fuente, arreglo sincronizado e inversión; existe compromiso entre resolución y profundidad \\
MAM / ReMi / SPAC & Dispersión del ruido ambiental & Extiende la investigación a frecuencias bajas y mayores profundidades sin fuente controlada & Depende del campo de ruido, la geometría y registros más largos \\
HVSR & Frecuencia o período fundamental del sitio & Una estación triaxial y despliegue rápido & Por sí solo no determina un perfil $V_S$ único \\
ERT & Distribución de resistividad eléctrica & Sensible a saturación, arcillas, cavidades y cambios litológicos & La resistividad no es una medida directa de rigidez mecánica \\
\bottomrule
\end{tabularx}
\end{table}

\subsubsection{Los métodos clásicos siguen siendo necesarios}

SPT y CPT son valiosos porque entregan información de contacto y se encuentran
integrados en la práctica geotécnica. Los ensayos downhole y crosshole, por su parte,
permiten medir velocidades sísmicas dentro o entre perforaciones. Su principal
desventaja no es la calidad del dato, sino la representatividad espacial y la logística:
un sondeo describe muy bien su entorno inmediato, pero la extrapolación entre sondeos
puede ocultar variaciones laterales.

Por eso, la pregunta correcta no es ``¿sondeo o geofísica?'', sino ``¿cómo usar los
sondeos para anclar una medición que cubra un volumen mayor?''. El caso de CentrePort
es ilustrativo: los datos CPT y de perforación se utilizaron para restringir la
parametrización de la inversión de ondas superficiales~\cite{Vantassel2018}.

\subsubsection{Los métodos no invasivos intercambian contacto directo por cobertura}

La refracción, la resistividad y las ondas superficiales permiten estudiar líneas o
superficies sin perforar cada punto. Esa cobertura es útil para localizar cambios y
decidir dónde conviene realizar una verificación directa. El precio metodológico es que
el resultado se obtiene mediante un problema inverso: las propiedades se infieren a
partir de la respuesta observada y varias estructuras pueden explicar datos similares.

En consecuencia, la geofísica no debe presentarse como una ``radiografía'' perfecta.
Su utilidad está en incorporar información espacial y dinámica que los ensayos puntuales
no aportan con la misma facilidad.

\subsection{Por qué las ondas superficiales y por qué MASW}

Las ondas Rayleigh son dispersivas en un medio estratificado: las componentes de
distinta longitud de onda muestrean volúmenes diferentes del subsuelo y se propagan con
velocidades de fase distintas. Esa dependencia de velocidad con frecuencia contiene
información sobre la distribución de $V_S$.

El procedimiento MASW activo separa el problema en tres etapas:

\begin{enumerate}
    \item \textbf{Adquisición}: una fuente controlada genera ondas y un arreglo de
    receptores registra simultáneamente el campo de ondas.
    \item \textbf{Análisis de dispersión}: las señales se transforman al dominio
    frecuencia--velocidad para identificar una o varias curvas modales.
    \item \textbf{Inversión}: se buscan perfiles estratificados cuya dispersión teórica
    sea compatible con la observada.
\end{enumerate}

MASW evolucionó respecto de SASW al pasar de pares de receptores a un registro
multicanal simultáneo. El muestreo espacial facilita la separación de energía coherente,
reduce ambigüedades de fase y aumenta la productividad de campo~\cite{Park1999}. Los
métodos pasivos MAM, ReMi o SPAC no compiten necesariamente con el activo: pueden
extender la curva hacia frecuencias más bajas y, por tanto, mayores profundidades.

La elección de MASW activo para este proyecto responde a que permite generar una
excitación repetible, trabajar con una geometría conocida y evaluar de manera directa
la cadena instrumental multicanal desarrollada. Además, el mismo sistema puede explorar
en el futuro adquisiciones pasivas mediante registros continuos.

La elección no elimina sus limitaciones:

\begin{itemize}
    \item la profundidad depende de las longitudes de onda realmente observadas con
    relación señal--ruido suficiente;
    \item la apertura, el espaciamiento y el \emph{offset} del arreglo condicionan la
    banda válida;
    \item los modos superiores pueden confundirse con el modo fundamental;
    \item la inversión es no única y debe reportarse con incertidumbre;
    \item errores de fase entre canales pueden convertirse directamente en errores de
    velocidad de fase.
\end{itemize}

El proyecto InterPACIFIC y las guías derivadas demostraron que los métodos de ondas
superficiales pueden producir resultados reproducibles entre equipos y analistas cuando
se documentan la geometría, la banda válida, el procesamiento y la incertidumbre. A la
vez, mostraron que las diferencias del perfil final dependen fuertemente de la
parametrización y la inversión~\cite{Foti2018}. Esto acota el papel del instrumento: debe
producir datos sincronizados y trazables de calidad conocida, pero no puede eliminar por
sí solo la no unicidad del subsuelo.

\subsection{De la señal medida al modelo \texorpdfstring{$V_S(z)$}{Vs(z)}: física y ecuaciones}

La cadena que suele resumirse como

\begin{center}
$V_{R,\mathrm{aparente}}\longrightarrow V_S\longrightarrow\text{modelo}$
\end{center}

contiene en realidad dos problemas distintos. Primero se estima, a partir de las
señales multicanal, cómo cambia la velocidad de fase de las ondas Rayleigh con la
frecuencia. Después se buscan las propiedades de un medio estratificado que puedan
producir esa dispersión. Esta distinción es esencial: en un suelo formado por capas no
se mide directamente la $V_S$ de una profundidad determinada.

\subsubsection{Del registro espacio--tiempo a la velocidad de fase aparente}

Sea $u(x_j,t)$ el desplazamiento o, en la práctica, la señal proporcional a velocidad
registrada por el geófono ubicado en $x_j$. Para cada canal se calcula la transformada
temporal de Fourier

\begin{equation}
    U(x_j,\omega)=\int_{-\infty}^{\infty}u(x_j,t)
    e^{i\omega t}\,\mathrm{d}t
    =A(x_j,\omega)e^{i\phi(x_j,\omega)}.
\end{equation}

Si una componente de frecuencia $f=\omega/(2\pi)$ se propaga como una onda plana, su
fase cambia aproximadamente como $\phi(x,\omega)=\phi_0-kx$, con número de onda
$k=\omega/c_R$. Entre dos receptores separados una distancia $\Delta x$ se obtendría
entonces

\begin{equation}
    c_{R,\mathrm{aparente}}(f)
    =\frac{\omega\,\Delta x}{|\Delta\phi(f)|}
    =\frac{2\pi f\,\Delta x}{|\Delta\phi(f)|}.
    \label{eq:velocidad-fase-dos-canales}
\end{equation}

El problema práctico de esta expresión es que la fase está definida módulo $2\pi$,
puede contaminarse con ruido y puede contener simultáneamente varios modos. MASW utiliza
todos los canales para reducir esa ambigüedad. En el método de desplazamiento de fase
se prueba una velocidad $c$ y se corrige la fase esperada en cada posición. Una forma
normalizada del espectro es~\cite{Park1998}

\begin{equation}
    S(\omega,c)=\frac{1}{N}\left|
    \sum_{j=1}^{N}
    \frac{U(x_j,\omega)}{|U(x_j,\omega)|}
    \exp\!\left(i\frac{\omega x_j}{c}\right)
    \right|.
    \label{eq:phase-shift}
\end{equation}

Cuando la velocidad ensayada coincide con la velocidad de fase de una onda presente,
las contribuciones de los receptores se alinean y $S$ presenta un máximo. Al seleccionar
esos máximos para cada frecuencia se obtiene la curva observada

\begin{equation}
    c_R^{\mathrm{obs}}(f)=\operatorname*{arg\,max}_{c}S(2\pi f,c).
\end{equation}

Se la denomina velocidad de fase \emph{aparente} porque se infiere desde un arreglo
finito y puede incluir efectos de ruido, modos, geometría y variación lateral. Bajo la
aproximación de estratos horizontales representa la dispersión modal de Rayleigh que se
utilizará en la inversión.

\subsubsection{Frecuencia, longitud de onda y profundidad}

Cada punto de la curva tiene asociada una longitud de onda

\begin{equation}
    \lambda_R(f)=\frac{c_R(f)}{f}.
    \label{eq:longitud-onda}
\end{equation}

Las longitudes de onda cortas concentran mayor sensibilidad cerca de la superficie;
las largas involucran un volumen más profundo. No existe, sin embargo, una igualdad
universal del tipo $z=\lambda/2$: cada medición posee una función de sensibilidad
distribuida y solapada con la de otras frecuencias. La regla de una fracción de la
longitud de onda sirve para diseñar e interpretar preliminarmente el ensayo, no para
asignar cada velocidad a una profundidad exacta.

Por ejemplo, un máximo de dispersión de \SI{240}{\meter\per\second} a
\SI{20}{\hertz} corresponde a $\lambda_R=\SI{12}{\meter}$. Ese dato restringe la
combinación de propiedades de las capas someras sensibles a esa longitud de onda; no
significa que $V_S$ a \SI{6}{\meter} sea \SI{240}{\meter\per\second}. Un punto de
\SI{500}{\meter\per\second} a \SI{5}{\hertz}, en cambio, tiene
$\lambda_R=\SI{100}{\meter}$ y aporta información integrada de un volumen mucho más
profundo.

\subsubsection{Relación física entre \texorpdfstring{$c_R$ y $V_S$}{cR y Vs}}

En un semiespacio elástico, homogéneo e isótropo, la velocidad de Rayleigh se relaciona
con $V_S$ y $V_P$. Si $\xi=c_R/V_S$ y $\kappa=(V_S/V_P)^2$, la raíz física satisface

\begin{equation}
    \xi^6-8\xi^4+(24-16\kappa)\xi^2-16(1-\kappa)=0.
    \label{eq:rayleigh-homogeneo}
\end{equation}

Para una relación de Poisson $\nu=0{,}25$, $\kappa=1/3$ y
$c_R\simeq0{,}9194V_S$~\cite{Xia1999}. Sólo bajo ese modelo homogéneo podría utilizarse
la aproximación $V_S\simeq c_R/0{,}9194$. Así, para el ejemplo anterior se obtendría
$V_S\simeq\SI{261}{\meter\per\second}$, pero sería la velocidad de un semiespacio
equivalente, no el valor de una capa a una profundidad particular.

En un terreno estratificado, las condiciones de continuidad de esfuerzos y
desplazamientos en cada interfaz producen una ecuación característica cuyos ceros
dependen de la frecuencia y del modelo completo:

\begin{equation}
    F\!\left(f,c_R;\boldsymbol{V}_S,\boldsymbol{V}_P,
    \boldsymbol{\rho},\boldsymbol{h}\right)=0.
    \label{eq:forward-rayleigh}
\end{equation}

Aquí $\boldsymbol{h}$ contiene los espesores y $\boldsymbol{V}_S$,
$\boldsymbol{V}_P$ y $\boldsymbol{\rho}$ las propiedades de cada estrato. Para cada
frecuencia, las raíces de la ecuación~\eqref{eq:forward-rayleigh} son las velocidades
de fase teóricas de los distintos modos. En el rango habitual de ingeniería la
dispersión de Rayleigh es especialmente sensible a $V_S$ y, después, a los espesores;
por eso $V_P$ y $\rho$ suelen fijarse o restringirse con información auxiliar
~\cite{Xia1999}.

\subsubsection{De la curva de dispersión a una familia de modelos}

La inversión parte de los datos
$\boldsymbol{d}=[c_R^{\mathrm{obs}}(f_1),\ldots,c_R^{\mathrm{obs}}(f_M)]$ y propone un
modelo

\begin{equation}
\begin{split}
    \boldsymbol{m}=[&h_1,\ldots,h_{L-1},
    V_{S,1},\ldots,V_{S,L},\\
    &V_{P,1},\ldots,V_{P,L},
    \rho_1,\ldots,\rho_L].
\end{split}
\end{equation}

Un solucionador directo calcula
$\boldsymbol{g}(\boldsymbol{m})=\boldsymbol{c}_R^{\mathrm{calc}}$ mediante la
ecuación~\eqref{eq:forward-rayleigh}. La inversión modifica el modelo para minimizar,
por ejemplo,

\begin{equation}
    \Phi(\boldsymbol{m})=
    \left\|\boldsymbol{W}_d
    [\boldsymbol{d}-\boldsymbol{g}(\boldsymbol{m})]\right\|_2^2
    +\lambda\,R(\boldsymbol{m}),
    \label{eq:objetivo-inversion}
\end{equation}

donde $\boldsymbol{W}_d$ pondera cada dato según su incertidumbre y
$R(\boldsymbol{m})$ penaliza modelos incompatibles con restricciones elegidas, por
ejemplo variaciones excesivamente abruptas. En una inversión linealizada, la corrección
puede expresarse como

\begin{equation}
    (\boldsymbol{J}^{T}\boldsymbol{W}\boldsymbol{J}+\lambda\boldsymbol{L}^{T}
    \boldsymbol{L})\,\Delta\boldsymbol{m}
    =\boldsymbol{J}^{T}\boldsymbol{W}
    [\boldsymbol{d}-\boldsymbol{g}(\boldsymbol{m})],
\end{equation}

con $J_{ij}=\partial c_R(f_i)/\partial m_j$. El proceso se repite con
$\boldsymbol{m}_{k+1}=\boldsymbol{m}_k+\Delta\boldsymbol{m}$, o se reemplaza por una
búsqueda global; aquí $\boldsymbol{W}=\boldsymbol{W}_d^T\boldsymbol{W}_d$. En ambos
casos la salida honesta no es una única ``fotografía'', sino
un conjunto de modelos aceptables y sus bandas de variación.

Una vez estimado el perfil $V_S(z)$ pueden calcularse magnitudes de ingeniería. Para
capas que completan los primeros \SI{30}{\meter},

\begin{equation}
    V_{S,30}=\frac{30}{\displaystyle\sum_i h_i/V_{S,i}},
    \qquad
    G_{\max,i}=\rho_i V_{S,i}^{2}.
\end{equation}

La cadena física completa queda, por tanto,

\begin{center}
\small
$u_j(t)\rightarrow U_j(f)\rightarrow S(f,c)\rightarrow
c_R^{\mathrm{obs}}(f)\rightarrow
\text{inversión estratificada}\rightarrow
\{h_i,V_{S,i},V_{P,i},\rho_i\}\rightarrow V_S(z),V_{S,30},G_{\max}$.
\end{center}

Esto también explica por qué la instrumentación importa: error de fase entre canales,
una banda recortada por el transductor o una geometría mal documentada alteran la curva
de dispersión antes de que el algoritmo de inversión pueda intervenir.

\subsection{Soluciones comerciales: evidencia de interés industrial}

La existencia de varias familias de sismógrafos profesionales confirma que la
adquisición sísmica somera no es un problema puramente académico. Los fabricantes han
invertido en mejorar exactamente los aspectos que aparecen en campo: resolución y rango
dinámico, sincronización, control sencillo, autonomía, escalabilidad y reducción de
cableado.

Por eso no resulta correcto presentar esta parte como ``problemas de los MASW
comerciales''. Son soluciones maduras y, en varios casos, muy capaces. Lo que interesa
es observar qué arquitectura adopta cada una y qué necesidad industrial demuestra.

\begin{table}[ht]
\centering
\footnotesize
\caption{Soluciones comerciales y señal de demanda industrial. Las prestaciones provienen de las páginas oficiales de los fabricantes.}
\begin{tabularx}{\textwidth}{@{}>{\raggedright\arraybackslash}p{0.16\textwidth} >{\raggedright\arraybackslash}p{0.22\textwidth} Y Y@{}}
\toprule
\textbf{Producto} & \textbf{Arquitectura} & \textbf{Qué resuelve} & \textbf{Qué demuestra} \\
\midrule
Geometrics GeodeFlex & 24 canales analógicos centralizados; control web; Wi-Fi en versión Core; GPS integrado & Adquisición simultánea de 24 bits, operación desde PC, tableta o teléfono y registro autónomo & La interfaz web y la sincronización forman parte del valor industrial, no son accesorios \\
PASI GEA24 & Sismógrafo compacto de 24 canales conectado por USB a una computadora; ampliable a 48 & Portabilidad, alimentación USB y uso multipropósito: MASW, Vs30, refracción, tomografía y pozo & Existe demanda por una plataforma compacta que reutilice la misma adquisición en varios métodos \\
DMT SUMMIT X One & Unidades remotas distribuidas sobre un cable ligero de dos conductores para datos y energía & Geometrías flexibles, despliegue rápido y escalabilidad desde más de 24 hasta miles de canales & Reducir el cableado y distribuir la adquisición tiene valor operativo real \\
Geometrics Atom-3C & Unidades triaxiales autónomas, temporizadas por GPS, sin cable de tendido; descarga por Wi-Fi & Adquisición pasiva prolongada sin fuente, computadora ni interconexión entre estaciones & La autonomía y la eliminación de cables son suficientemente valiosas para justificar una familia específica de producto \\
\bottomrule
\end{tabularx}
\end{table}

\subsubsection{Geometrics GeodeFlex: control moderno sobre una adquisición centralizada}

GeodeFlex integra 24 canales con conversión sigma--delta de 24 bits y muestreo
simultáneo de hasta \SI{32}{\kilo\hertz} por canal. La versión Core incorpora
controlador Linux, GPS, Wi-Fi y una aplicación web operable desde computadora, tableta o
teléfono. El sistema puede realizar registro autónomo sincronizado por GPS para ensayos
pasivos y de monitoreo~\cite{GeometricsGeodeFlex}.

Es importante interpretar bien la palabra ``inalámbrico'': el Wi-Fi simplifica el
control, pero los 24 geófonos analógicos continúan conectados mediante el cable de
tendido y el disparo activo es cableado. El producto demuestra que una interfaz web
mejora la operación de campo incluso cuando la arquitectura de adquisición sigue siendo
centralizada.

\subsubsection{PASI GEA24: compactación del sistema profesional clásico}

GEA24 integra 24 canales, conversión de 24 bits, interfaz USB y alimentación desde la
computadora. Dos unidades pueden sincronizarse para obtener 48 canales. El fabricante lo
orienta a refracción, tomografía, downhole, crosshole, $V_{S,30}$, MASW y reflexión
somera~\cite{PASIGEA24}. El caso de Delhi confirma que no se trata de una lista teórica
de aplicaciones: el equipo fue utilizado efectivamente para caracterizar seis sectores
de un campus~\cite{Gautam2020}.

Su arquitectura conserva los cables sísmicos de 12 o 24 canales y depende de una
computadora externa, pero reduce el volumen del sismógrafo y elimina una batería de
\SI{12}{\volt} independiente. Representa una optimización clara del esquema centralizado.

\subsubsection{DMT SUMMIT X One: adquisición distribuida sin abandonar el cable}

SUMMIT X One desplaza la conversión hacia unidades remotas que pueden conectarse en
posiciones flexibles sobre una línea de dos conductores. Ese cable distribuye tanto
energía como datos. El fabricante ofrece configuraciones desde más de 24 canales para
ingeniería hasta despliegues 2D/3D de más de 3000 canales, además de registro continuo
para aplicaciones pasivas~\cite{DMTSummitXOne}.

La lección de arquitectura es relevante: distribuir nodos no obliga a eliminar todo
cable. Un enlace físico ligero puede simplificar simultáneamente alimentación,
telemetría y sincronización. El producto demuestra que la logística del arreglo es un
problema industrial suficientemente importante como para rediseñar el sistema completo
alrededor de ella.

\subsubsection{Geometrics Atom-3C: nodos autónomos para adquisición pasiva}

Atom-3C elimina el cable de tendido mediante unidades autónomas de tres canales. Cada
unidad incorpora conversión de 24 bits efectiva, temporización GPS, almacenamiento
interno, batería y descarga por Wi-Fi. El fabricante declara hasta 70 horas de operación
por carga y lo orienta a microtremores, clasificación $V_{S,30}$, estructuras profundas
y microzonificación~\cite{GeometricsAtom3C}.

Atom-3C resuelve muy bien la adquisición pasiva, pero no es un sistema MASW activo con
disparo de fuente. Esa diferencia evita una comparación engañosa: quitar los cables es
más sencillo cuando cada nodo registra continuamente con reloj GPS que cuando todos
deben responder de forma determinista a una fuente impulsiva.

\subsection{Precios de referencia: dimensión económica de la brecha}

Los fabricantes venden configuraciones diferentes y varios trabajan exclusivamente
mediante cotización. Por eso, un precio aislado puede ser engañoso: debe indicarse el
año, la cantidad de canales y si incluye geófonos, cables, fuente, software, transporte
e impuestos. La tabla siguiente utiliza listas públicas encontradas, conserva la moneda
original y no ajusta inflación ni tipo de cambio.

\begin{table}[ht]
\centering
\scriptsize
\caption{Precios públicos de referencia. Son antecedentes fechados, no cotizaciones vigentes.}
\begin{tabularx}{\textwidth}{@{}>{\raggedright\arraybackslash}p{0.15\textwidth} >{\raggedright\arraybackslash}p{0.22\textwidth} Y Y@{}}
\toprule
\textbf{Sistema} & \textbf{Referencia pública} & \textbf{Alcance declarado} & \textbf{Lectura correcta} \\
\midrule
PASI GEA24 & 4.125~EUR sólo instrumento; 8.525~EUR sistema MASW de 24 canales (enero de 2022) & Consola, dos cables de 12 canales, 24 geófonos verticales de \SI{4.5}{\hertz}, disparo de maza, cable de disparo y placa; excluye PC, software de análisis, transporte e impuestos & Es la comparación más cercana a un arreglo MASW activo completo, pero es precio histórico \emph{ex works} Turín~\cite{PASIPrice2022} \\
Geometrics & 18.700~USD por un Geode G24 de 24 canales (lista 2023) & Sismógrafo, interfaz digital, accesorios de alimentación/disparo y SeisImager/2D Lite; sin geófonos ni cable sísmico & Sirve como orden de magnitud de la generación anterior; \textbf{no es el precio de GeodeFlex}, que actualmente se ofrece por cotización~\cite{GeometricsPrice2023,GeometricsGeodeFlex} \\
DMT SUMMIT X One & Precio vigente sólo por solicitud & La arquitectura puede variar en número de unidades remotas, recolectores, tramos de cable, disparos y accesorios & La ausencia de una cifra pública refleja que se cotiza un sistema configurado, no una caja estándar~\cite{GeoDeviceSummitPrice} \\
Geometrics Atom & 1.800~USD por unidad de adquisición y 1.160~USD por geófono triaxial de \SI{2}{\hertz} (lista 2023) & Doce estaciones sumarían 35.520~USD; con un punto de acceso de 900~USD y dos cargadores de seis puertos de 430~USD, 37.280~USD antes de software, transporte e impuestos & Es una estimación aritmética para doce estaciones pasivas a partir de componentes de catálogo, no una cotización Atom-3C actual~\cite{GeometricsPrice2023} \\
MicDAC académico & Aproximadamente 255~EUR (2020) & Hardware completo con sensores para tres componentes, 12 bits y \SI{200}{\hertz}, destinado a HVSR & Demuestra acceso económico, pero no equivale a 24 canales sincronizados ni a un sistema MASW activo~\cite{Kafadar2020} \\
\bottomrule
\end{tabularx}
\end{table}

La comparación no pretende concluir que dos sistemas son equivalentes por costar menos.
El GEA24 completo ofrece una referencia histórica del orden de miles de euros; el Geode
de 24 canales, aun sin sensores ni tendido, aparece en decenas de miles de dólares; y
una red pasiva Atom escala aproximadamente con cada estación autónoma. En contraste,
MicDAC muestra que un registrador académico de pocos canales puede construirse por
cientos de euros, pero todavía no resuelve sincronización, disparo activo, cantidad de
canales, robustez ni flujo de procesamiento.

\subsection{La brecha no es tecnológica: es de acceso e integración económica}

Los cuatro productos anteriores prueban que la industria ya resolvió, con distintas
arquitecturas, buena parte de los problemas de adquisición. La brecha que motiva esta
tesis debe formularse con más precisión:

\begin{quote}
\textbf{Existe instrumentación profesional capaz, pero falta una plataforma económica,
reproducible y modificable que integre adquisición multicanal sincronizada, operación
de campo y procesamiento para investigación local.}
\end{quote}

La dificultad de acceso no se reduce al precio publicado del sismógrafo: también incluye
sensores, cables, disparo, software, importación, impuestos, soporte y reparación. Para
una empresa especializada esa inversión puede justificarse por productividad y
confiabilidad. Para un laboratorio pequeño, una universidad o una etapa exploratoria,
la barrera es diferente.

La literatura académica confirma esa demanda por alternativas. Kafadar desarrolló
MicDAC, un sistema abierto de tres componentes basado en Arduino para mediciones HVSR,
con una frecuencia de muestreo de \SI{200}{\hertz}, conversión de 12 bits y un costo de
componentes, incluidos sensores, estimado en 255~EUR en 2020~\cite{Kafadar2020}.
Geophonino demostró que un registrador abierto y económico de un canal puede adquirir
señales de un geófono y compararse favorablemente con equipos comerciales en su alcance
de prueba~\cite{Llorens2016}.

Estos antecedentes son valiosos porque muestran que el acondicionamiento y la
digitalización económicos son posibles. Al mismo tiempo, dejan visible el salto que
falta para MASW activo:

\begin{itemize}
    \item pasar de uno o tres canales a un arreglo de nodos sincronizados;
    \item coordinar una fuente impulsiva y conservar una referencia temporal común;
    \item controlar ganancias, calibración, almacenamiento y estado de cada nodo;
    \item operar el sistema en campo sin depender de infraestructura externa;
    \item preservar datos crudos y metadatos de geometría;
    \item integrar dispersión e inversión sin ocultar las decisiones del operador.
\end{itemize}

\begin{table}[ht]
\centering
\small
\caption{Posicionamiento de la brecha abordada por la tesis.}
\begin{tabularx}{\textwidth}{@{}>{\raggedright\arraybackslash}p{0.22\textwidth} Y Y Y@{}}
\toprule
\textbf{Necesidad} & \textbf{Equipo profesional} & \textbf{Prototipo económico típico} & \textbf{Objetivo de esta tesis} \\
\midrule
Adquisición multicanal & Resuelta y validada & Uno o pocos canales & Nodos ampliables con captura coordinada \\
Sincronización & GPS, cable o electrónica dedicada & Frecuentemente no necesaria en un solo nodo & Sincronización distribuida medida y documentada \\
Operación de campo & Software del fabricante y gabinete robusto & Interfaz de laboratorio o computadora conectada & Interfaz web local sin dependencia de Internet \\
Modificabilidad & Producto terminado y soportado & Hardware abierto pero alcance limitado & Cadena completa inspeccionable y reconfigurable \\
Procesamiento & Ecosistema comercial integrado & Análisis específico o externo & Datos crudos, ingesta, revisión, dispersión e inversión \\
\bottomrule
\end{tabularx}
\end{table}

El objetivo no es afirmar que un prototipo académico sustituye a GeodeFlex, GEA24,
SUMMIT X One o Atom-3C en confiabilidad, robustez o soporte. El aporte consiste en
construir y evaluar una plataforma accesible que permita estudiar las decisiones que
esos productos encapsulan: transducción, acondicionamiento, digitalización,
sincronización, transmisión, interfaz y procesamiento.

\subsection{Acotamiento del proyecto}

La revisión anterior conduce a una cadena funcional concreta:

\begin{center}
\small
Generación de señal $\longrightarrow$ transductor $\longrightarrow$
acondicionamiento e instrumentación $\longrightarrow$ digitalización

$\longrightarrow$ transmisión e interfaz de campo $\longrightarrow$ ingesta
$\longrightarrow$ procesamiento e inversión con interfaz de análisis.
\end{center}

Cada bloque presenta una pregunta verificable. El transductor define la banda que puede
observarse; el acondicionamiento determina ruido, ganancia y rango dinámico; la
digitalización y sincronización condicionan la coherencia entre canales; la transmisión
define cómo se despliega el arreglo; las interfaces determinan si el sistema puede
operarse y auditarse; y el procesamiento convierte los registros en un modelo con
incertidumbre.

La pregunta que abre el resto de la presentación puede formularse así:

\begin{quote}
\textbf{¿Cómo integrar, con componentes accesibles, una cadena multicanal de
adquisición y análisis para ondas superficiales que sea operable en campo, conserve la
trazabilidad de los datos y permita evaluar honestamente sus limitaciones?}
\end{quote}

Las secciones siguientes estudiarán cada bloque de esa cadena: selección y banda del
transductor, evolución del acondicionamiento, digitalización, comunicación distribuida,
interfaces de operación e ingesta, y finalmente procesamiento e inversión.
